% Combined TeX source bundle
% Title: Omar Khayyam, Treatise on Algebra
% Note: constituent files are preserved below in sources/. Some are standalone
% documents with their own preambles; this combined file is for inspection,
% comparison, and source review rather than guaranteed one-command compilation.


% ============================================================
% Source: translations/en/islamic/omar-khayyam-sinaat-al-jabr_bilingual.tex
% ============================================================

%% ========================================================================
%% Omar Khayyam - Sina'at al-Jabr wa-al-Muqabala
%% Bilingual English-Arabic Edition
%% Left column: Arabic text  |  Right column: English translation
%% Based on Kasir 1931 translation style
%% ========================================================================
\documentclass[11pt,a4paper]{book}

%% -------- Core Packages --------
\usepackage{fontspec}
\usepackage{polyglossia}
\usepackage{amsmath,amssymb}
\usepackage{geometry}
\usepackage{graphicx}
\usepackage{xcolor}
\usepackage{titlesec}
\usepackage{fancyhdr}
\usepackage{enumitem}
\usepackage{array,booktabs,longtable,multirow}
\usepackage{hyperref}
\usepackage{microtype}
\usepackage{tocloft}
\usepackage{tikz}

\usetikzlibrary{arrows.meta,calc,intersections,patterns}

%% -------- Page Geometry --------
\geometry{
  paperwidth=210mm,paperheight=297mm,
  left=18mm,right=18mm,top=25mm,bottom=25mm,
  headheight=14pt,footskip=30pt
}

%% -------- Language Setup --------
\setmainlanguage{english}
\setotherlanguage{arabic}

%% -------- Font Configuration --------
\setmainfont{lmroman10-regular.otf}[
  Path=fonts/latin-modern/,
  BoldFont=lmroman10-bold.otf,
  ItalicFont=lmroman10-italic.otf,
  BoldItalicFont=lmroman10-bolditalic.otf,
  SlantedFont=lmromanslant10-regular.otf,
  Renderer=Harfbuzz
]
\setsansfont{lmsans10-regular.otf}[
  Path=fonts/latin-modern/,
  BoldFont=lmsans10-bold.otf,
  ItalicFont=lmsans10-oblique.otf,
  Renderer=Harfbuzz
]
\setmonofont{lmmono10-regular.otf}[
  Path=fonts/latin-modern/,
  Scale=MatchLowercase
]

%% Arabic Font
\TeXXeTstate=1
\newfontfamily\arabicfont[Script=Arabic,Scale=1.15,Path=/usr/share/fonts/truetype/noto/]{NotoNaskhArabic-Regular.ttf}
\newfontfamily\arabicfontsf[Script=Arabic,Scale=1.1]{Noto Sans Arabic}

%% -------- Colors --------
\definecolor{titlecolor}{RGB}{45,55,72}
\definecolor{sectioncolor}{RGB}{55,65,81}
\definecolor{notecolor}{RGB}{120,53,15}
\definecolor{rulecolor}{RGB}{180,160,140}
\definecolor{arabicbg}{RGB}{252,250,245}
\definecolor{translationgray}{RGB}{80,80,80}
\definecolor{footnotegreen}{RGB}{0,80,40}
\definecolor{propline}{RGB}{139,90,43}
\definecolor{propbg}{RGB}{255,254,252}

%% -------- Section Formatting --------
\titleformat{\chapter}[display]
  {\normalfont\huge\bfseries\color{titlecolor}}
  {\chaptertitlename\ \thechapter}{20pt}{\Huge}
\titleformat{\section}
  {\normalfont\Large\bfseries\color{sectioncolor}}
  {\thesection}{1em}{}
\titleformat{\subsection}
  {\normalfont\large\bfseries\color{sectioncolor}}
  {\thesubsection}{1em}{}

%% -------- Header/Footer --------
\pagestyle{fancy}
\fancyhf{}
\fancyhead[L]{\small\textit{Omar Khayyam -- Sina'at al-Jabr (Bilingual Edition)}}
\fancyhead[R]{\small\thepage}
\renewcommand{\headrulewidth}{0.4pt}
\renewcommand{\footrulewidth}{0pt}

%% -------- Bilingual Column Macros --------
\newlength{\arabiccolwidth}
\setlength{\arabiccolwidth}{0.47\textwidth}
\newlength{\englishcolwidth}
\setlength{\englishcolwidth}{0.47\textwidth}

%% Arabic left column (RTL)
\newcommand{\arabiccol}[1]{%
  \begin{minipage}[t]{\arabiccolwidth}%
  \raggedleft\arabicfont\beginR #1\endR%
  \end{minipage}%
}

%% English right column
\newcommand{\englishcol}[1]{%
  \begin{minipage}[t]{\englishcolwidth}%
  \raggedright\color{translationgray}#1%
  \end{minipage}%
}

%% Separator
\newcommand{\colseparator}{\hfill\vrule width 0.4pt\hfill}

%% Full bilingual line: Arabic | English
\newcommand{\bilingualline}[2]{%
  \noindent\arabiccol{#1}\colseparator\englishcol{#2}\par\vspace{0.8em}%
}

%% -------- Custom Commands --------
\newcommand{\longline}{\noindent\rule{\textwidth}{0.4pt}}
\newcommand{\shortline}{\noindent\rule{0.5\textwidth}{0.4pt}\par}
\newcommand{\translatornote}[1]{\footnote{\textcolor{footnotegreen}{\textbf{[Tr.]} #1}}}

%% Inline Arabic
\newcommand{\ar}[1]{{\arabicfont\beginR #1\endR}}

%% Proposition box using simple rules
\newenvironment{propositionbox}[1][]{%
  \par\vspace{0.5em}%
  \noindent\rule{\textwidth}{0.6pt}\par\vspace{0.2em}%
  \ifx&#1&\else\noindent\textbf{#1}\par\vspace{0.3em}\fi%
}{%
  \par\vspace{0.2em}\noindent\rule{\textwidth}{0.6pt}\par\vspace{0.5em}%
}

%% Definition bar using simple rules
\newenvironment{definitionbar}{%
  \par\vspace{0.5em}%
  \noindent\rule{\textwidth}{0.4pt}\par\vspace{0.2em}%
}{%
  \par\vspace{0.2em}\noindent\rule{\textwidth}{0.4pt}\par\vspace{0.5em}%
}

%% -------- Hyperref Setup --------
\hypersetup{
  colorlinks=true,
  linkcolor=sectioncolor,
  citecolor=sectioncolor,
  urlcolor=sectioncolor,
  pdfencoding=auto,
  pdftitle={Omar Khayyam - Sina'at al-Jabr (Bilingual Edition)},
  pdfauthor={Omar Khayyam},
  pdfsubject={Bilingual Arabic-English Edition of Treatise on Algebra}
}

%% -------- TOC Depth --------
\setcounter{tocdepth}{2}
\setcounter{secnumdepth}{2}

%% ========================================================================
\begin{document}
%% ========================================================================

%% ---- Title Page --------------------------------------------------------
\begin{titlepage}
\begin{center}
\vspace*{1.5cm}

{\color{titlecolor}\rule{\textwidth}{1.5pt}}
\vspace{1cm}

{\Huge\bfseries\color{titlecolor} ʿUmar al-Khayyām}
\vspace{0.5cm}

{\Large\itshape (1048--1131 CE / 430--517 AH)}

\vspace{1.5cm}

{\color{titlecolor}\rule{0.7\textwidth}{0.8pt}}

\vspace{1.5cm}

{\arabicfont\LARGE\beginR صناعة الجبر والمقابلة\endR}

\vspace{0.3cm}

{\Large\itshape Sinaʿat al-Jabr wa-al-Muqābalah}

\vspace{0.5cm}

{\large\itshape Treatise on the Art of Algebra and Almucabala}

\vspace{1.5cm}

{\color{titlecolor}\rule{0.7\textwidth}{0.8pt}}

\vspace{1.5cm}

{\Large\bfseries\color{sectioncolor} Bilingual Edition}

\vspace{0.3cm}

{\large\color{translationgray}
Arabic Text \quad\textbar\quad English Translation}

\vspace{1.5cm}

{\large Translation in the style of}

{\large\itshape D. S. Kasir, \textit{The Algebra of Omar Khayyam} (1931)}

\vspace{2cm}

{\color{titlecolor}\rule{0.4\textwidth}{0.4pt}}

\vspace{0.5cm}

{\small The manuscript was copied in Lahore, India, in the 13th century.}

{\small It contains the earliest systematic treatment of cubic equations,}

{\small with all 25 types solved through geometric constructions using conic sections.}

\vspace{0.5cm}

{\small\textit{Bilingual typesetting via XeLaTeX, polyglossia, and Noto Naskh Arabic}}

\vspace{0.5cm}

{\footnotesize This edition prepared for scholarly and educational purposes.}

\end{center}
\end{titlepage}

%% ---- Copyright Page ----------------------------------------------------
\clearpage
\thispagestyle{empty}
\vspace*{8cm}
\begin{center}
{\small\textit{Sinaʿat al-Jabr wa-al-Muqābalah}}

{\small\textit{Bilingual Arabic-English Edition}}

\vspace{1cm}

{\footnotesize Original Arabic text based on the 13th-century Lahore manuscript}

{\footnotesize (Smith Oriental MS 45, Columbia University Rare Book and Manuscript Library)}

\vspace{0.5cm}

{\footnotesize English translation follows the style of D. S. Kasir,}

{\footnotesize \textit{The Algebra of Omar Khayyam}, New York: Teachers College Press, 1931.}

\vspace{1cm}

{\footnotesize Key terminology:}

{\footnotesize kaʿb (\ar{كعب}) = cube \quad|\quad māl (\ar{مال}) = square}

{\footnotesize jadhr (\ar{جذر}) = root \quad|\quad ziyāda (\ar{زيادة}) = augmentation}

{\footnotesize nāqiṣ (\ar{ناقص}) = diminution/defect}

\end{center}
\clearpage

%% ---- Table of Contents -------------------------------------------------
\tableofcontents
\clearpage

%% ========================================================================
\chapter{Introduction}
\label{ch:intro}
%% ========================================================================

Omar Khayyam (\ar{عمر الخيام}), full name Ghiyāth al-Dīn Abū al-Fatḥ ʿUmar ibn Ibrāhīm al-Khayyāmī al-Nīsābūrī, was one of the most brilliant mathematicians, astronomers, and poets of the medieval Islamic world. Born in Nishapur in 1048 CE, he spent much of his life in Samarqand and Isfahan, where he produced groundbreaking works in mathematics and astronomy.

His \textit{Sinaʿat al-Jabr wa-al-Muqābalah} (also cited as \textit{Risāla fī al-Barāhīn ʿalā Masāʾil al-Jabr wa-al-Muqābala} --- ``Treatise on the Proofs of the Problems of Algebra and Almucabala'') represents the culmination of early Islamic algebra. In this work, Khayyam:
\begin{itemize}[noitemsep]
    \item Provides rigorous definitions of the fundamental quantities of algebra
    \item Classifies all equations of degree up to three into \textbf{25 types}
    \item Solves 14 irreducible cubic equations using geometric constructions with conic sections
    \item Distinguishes between arithmetical and geometrical solutions
\end{itemize}

This treatise, composed around 1079 CE, profoundly influenced later mathematicians including Sharaf al-Dīn al-Ṭūsī, and through translations into European languages, contributed to the development of algebra in the West.

%% ========================================================================
\section{The Author's Preface}
\label{sec:preface}
%% ========================================================================

Khayyam opens his treatise with a preface in which he acknowledges the achievements and limitations of his predecessors. The following is the Arabic text of the preface, with parallel English translation.

\begin{propositionbox}[The Author's Introduction \ar{(مقدمة المؤلف)}]

%% ---- Basmala ----
\bilingualline{%
الحمد لله رب العالمين، والعاقبة للمتقين، ولا عدوان إلا على الظالمين، وصلى الله على سيدنا محمد وآله أجمعين.%
}{%
Praise be to God, Lord of the Worlds, and the good end is for the righteous, and there is no enmity save against the oppressors; and may God bless our master Muhammad and all his family.%
}

\vspace{0.5cm}

%% ---- Main preface text ----
\bilingualline{%
أما بعد، فإني قد سبقني في هذا الفن جماعة من أهل العلم، قدموا فيه أصناف المسائل، وبينوا فيه أصول المسائل، وبرهنوا على بعضها، ولم يتفق لهم برهان على البعض، ولم يفرقوا فيها بين ما يمكن أن يوجد بالقطوع المخروطية، وبين ما لا يمكن أن يوجد إلا بأساليب هندسية أخرى.%
}{%
Now then: in this art there have preceded me a group of men of knowledge, who presented in it the species of problems, and explained therein the foundations of problems, and proved some of them, but did not succeed in proving all of them; and they did not distinguish between what can be found by conic sections and what cannot be found except by other geometrical methods.%
}

\vspace{0.3cm}

\bilingualline{%
وقد أملت أن أجمع هذا العلم في رسالة، وأبين فيها أصناف المسائل على وجه التمييز، وأضع لكل صنف برهاناً صحيحاً، وذلك أمر لم يسبقني إليه أحد من المتقدمين.%
}{%
I have desired to gather this science into a treatise, and to explain therein the species of problems with proper distinction, and to set down for each species a correct proof --- and this is something that none of the predecessors has preceded me in.%
}

\end{propositionbox}

%% ========================================================================
\section{On the Nature of This Treatise}
\label{sec:nature}
%% ========================================================================

Khayyam defines algebra as ``the science of equations'' (\ar{علم المعادلات}). He distinguishes carefully between:
\begin{enumerate}[noitemsep]
    \item \textbf{Simple equations:} Those reducible to linear or quadratic forms, solved by arithmetical or Euclidean methods
    \item \textbf{Compound (irreducible) equations:} Cubic equations requiring conic section constructions
\end{enumerate}

\begin{propositionbox}[Definition of Algebra]

\bilingualline{%
فنقول: إن الجبر والمقابلة علم يُستخرج به المجهول من المعلوم المفروض، إذا كان بينهما نسبة تقتضي ذلك.%
}{%
We say: Algebra and Almucabala is a science by which the unknown is extracted from the known datum, when between them there is a relation requiring that.%
}

\vspace{0.3cm}

\bilingualline{%
وهذا العلم يتضمن الأعداد والجذور والأموال والكعاب.%
}{%
And this science comprises numbers, roots, squares, and cubes.\translatornote{The Arabic terms are: ʿadad (number), jadhr (root), māl (square, lit.\ ``possession'' or ``wealth''), and kaʿb (cube, lit.\ ``die'' or ``solid with six square faces'').}%
}

\end{propositionbox}

\clearpage
%% ========================================================================
\chapter{Definitions of the Fundamental Quantities}
\label{ch:definitions}
%% ========================================================================

Khayyam begins with careful definitions of the algebraic quantities, following and extending the tradition of al-Khwārizmī:

\begin{definitionbar}
\bilingualline{%
أصل هذا العلم: أن العدد المفرد هو ما لا ينسب إلى جذر، ولا إلى مال، ولا إلى كعب، ولا إلى ما يتركب منها.%
}{%
The foundation of this science is: The simple number is that which is not referred to a root, nor to a square, nor to a cube, nor to what is compounded from them.%
}
\end{definitionbar}

%% ========================================================================
\section{The Root \ar{(الجذر)}}
\label{sec:root}
%% ========================================================================

\bilingualline{%
والجذر هو كل شيء مضروب في نفسه من الواحد وما فوقه من الأعداد، وما دونه من الكسور. فإذا ضرب في نفسه فهو مال، وإذا ضرب المال في الجذر فهو كعب، وإذا ضرب الكعب في الجذر فهو مال مال، وإذا ضرب مال المال في الجذر فهو مال كعب، وإذا ضرب مال الكعب في الجذر فهو كعب كعب.%
}{%
The root is every quantity multiplied by itself, beginning from one and above in whole numbers, and below in fractions. If it is multiplied by itself, it is a square; and if the square is multiplied by the root, it is a cube; and if the cube is multiplied by the root, it is square-square; and if square-square is multiplied by the root, it is square-cube; and if square-cube is multiplied by the root, it is cube-cube.\translatornote{Khayyam follows the tradition of naming higher powers by composition: māl māl = $x^4$, māl kaʿb = $x^5$, kaʿb kaʿb = $x^6$.}%
}

%% ========================================================================
\section{The Square \ar{(المال)} and the Cube \ar{(الكعب)}}
\label{sec:square-cube}
%% ========================================================================

\bilingualline{%
والمال هو ما ينتج من ضرب الجذر في نفسه. والكعب هو ما ينتج من ضرب المال في الجذر. وما فوق الكعب فهو مركب من هذه الأجناس.%
}{%
The square is what results from multiplying the root by itself. The cube is what results from multiplying the square by the root. And what is above the cube is compounded from these species.\translatornote{The Arabic term \textit{māl} (``wealth'', ``possession'') refers to $x^2$ because the square of an unknown quantity was likened to money or property. \textit{Kaʿb} literally means ``die'' (as in dice), a solid cube.}%
}

%% ========================================================================
\section{Classification of Algebraic Terms}
\label{sec:terms}
%% ========================================================================

Khayyam systematically enumerates the simple and compound species:

\begin{center}
\begin{tabular}{lll}
\toprule
\textbf{Arabic Term} & \textbf{Transliteration} & \textbf{Modern Notation} \\
\midrule
\ar{جذر} & \textit{jadhr} & $x$ \\
\ar{مال} & \textit{māl} & $x^2$ \\
\ar{كعب} & \textit{kaʿb} & $x^3$ \\
\ar{مال مال} & \textit{māl māl} & $x^4$ \\
\ar{مال كعب} & \textit{māl kaʿb} & $x^5$ \\
\ar{كعب كعب} & \textit{kaʿb kaʿb} & $x^6$ \\
\bottomrule
\end{tabular}
\end{center}

\clearpage
%% ========================================================================
\chapter{Classification of Equations}
\label{ch:classification}
%% ========================================================================

This section forms the heart of Khayyam's treatise. He classifies all equations of degree up to three into \textbf{twenty-five types}. Following the conventions of his time, all coefficients are positive, so equations like $x^3 + cx = b x^2 + d$ and $x^3 + b x^2 = cx + d$ are considered distinct types.

%% ========================================================================
\section{Equations Solvable by Arithmetical or Euclidean Methods}
\label{sec:simple}
%% ========================================================================

The first group consists of linear and quadratic equations, solvable without conic sections:

%% ------------------------------------------------------------------------
\subsection{Simple Equations (Types 1--6)}
\label{subsec:simple-eqs}
%% ------------------------------------------------------------------------

These six types were known since al-Khwārizmī and are solved by elementary methods.

\begin{enumerate}[label=\arabic*.,noitemsep]
    \item \ar{جذور تساوي أعداداً} --- Roots equal numbers: $bx = c$
    \item \ar{أموال تساوي جذوراً} --- Squares equal roots: $x^2 = bx$
    \item \ar{أموال تساوي أعداداً} --- Squares equal numbers: $x^2 = c$
    \item \ar{أموال وجذور تساوي أعداداً} --- Squares and roots equal numbers: $x^2 + bx = c$
    \item \ar{أموال وأعداد تساوي جذوراً} --- Squares and numbers equal roots: $x^2 + c = bx$
    \item \ar{جذور وأعداد تساوي أموالاً} --- Roots and numbers equal squares: $bx + c = x^2$
\end{enumerate}

%% ------------------------------------------------------------------------
\subsection{Binomial Cubic Equations (Types 7--12)}
\label{subsec:binomial}
%% ------------------------------------------------------------------------

\bilingualline{%
وأما الكعاب والأموال والجذور التي لا تختلف فيها الأجناس، فإنها تنقسم إلى اثني عشر نوعاً: منها ما يحل بالجبر، ومنها ما يحل بالقطوع المخروطية.%
}{%
As for cubes, squares, and roots in which the species are not mixed, they divide into twelve types: some solved by algebra, and some by conic sections.%
}

\begin{enumerate}[label=\arabic*.,resume,noitemsep]
    \item \ar{كعب يساوي مالاً} --- Cube equals square: $x^3 = b x^2$
    \item \ar{كعب يساوي جذراً} --- Cube equals root: $x^3 = c x$
    \item \ar{كعب يساوي عدداً} --- Cube equals number: $x^3 = d$
    \item \ar{مال يساوي كعباً} --- Square equals cube: $b x^2 = x^3$
    \item \ar{جذر يساوي كعباً} --- Root equals cube: $c x = x^3$
    \item \ar{عدد يساوي كعباً} --- Number equals cube: $d = x^3$
\end{enumerate}

These binomial equations reduce to simpler forms by division.

%% ------------------------------------------------------------------------
\subsection{Trinomial Equations Reducible to Quadratics (Types 13--16)}
%% ------------------------------------------------------------------------

\begin{enumerate}[label=\arabic*.,resume,noitemsep]
    \item \ar{كعاب وأموال تساوي جذوراً} --- Cubes and squares equal roots: $x^3 + b x^2 = c x$
    \item \ar{كعاب وجذور تساوي أموالاً} --- Cubes and roots equal squares: $x^3 + c x = b x^2$
    \item \ar{أموال وجذور تساوي كعاباً} --- Squares and roots equal cubes: $b x^2 + c x = x^3$
    \item \ar{كعاب وأموال وأعداد} --- [Special reducible forms]
\end{enumerate}

\clearpage
%% ========================================================================
\section{Irreducible Cubic Equations (The 14 Species)}
\label{sec:irreducible}
%% ========================================================================

The crowning achievement of Khayyam's treatise is his systematic solution of the \textbf{14 irreducible cubic equations}. These require the intersection of conic sections for their geometric solution. We present them here with their Arabic names and modern notation:

\begin{propositionbox}[The 14 Irreducible Cubic Equations \ar{(الأصناف الأربعة عشر)}]

\textbf{Group I: Three-term equations with successive powers}

\begin{enumerate}[label=\arabic*.,resume,noitemsep]
    \item \ar{كعب ومال يساويان عدداً} --- Cube and square equal number:
    $$x^3 + b x^2 = d$$
    
    \item \ar{كعب وعدد يساويان مالاً} --- Cube and number equal square:
    $$x^3 + d = b x^2$$
    
    \item \ar{كعب يساوي مالاً وعدداً} --- Cube equals square and number:
    $$x^3 = b x^2 + d$$
\end{enumerate}

\vspace{0.5cm}

\textbf{Group II: Three-term equations with mixed powers}

\begin{enumerate}[label=\arabic*.,resume,noitemsep]
    \setcounter{enumi}{18}
    \item \ar{كعب وجذور يساويان عدداً} --- Cube and roots equal number:
    $$x^3 + c x = d$$
    
    \item \ar{كعب وعدد يساويان جذوراً} --- Cube and number equal roots:
    $$x^3 + d = c x$$
    
    \item \ar{كعب يساوي جذوراً وعدداً} --- Cube equals roots and number:
    $$x^3 = c x + d$$
\end{enumerate}

\vspace{0.5cm}

\textbf{Group III: Four-term equations --- three terms on one side}

\begin{enumerate}[label=\arabic*.,resume,noitemsep]
    \setcounter{enumi}{21}
    \item \ar{كعب وأموال يساويان جذوراً وعدداً} --- Cube and squares equal roots and number:
    $$x^3 + b x^2 = c x + d$$
    
    \item \ar{كعب وجذور يساويان أموالاً وعدداً} --- Cube and roots equal squares and number:
    $$x^3 + c x = b x^2 + d$$
    
    \item \ar{كعب وأعداد يساويان أموالاً وجذوراً} --- Cube and numbers equal squares and roots:
    $$x^3 + d = b x^2 + c x$$
\end{enumerate}

\vspace{0.5cm}

\textbf{Group IV: Four-term equations --- two terms on each side}

\begin{enumerate}[label=\arabic*.,resume,noitemsep]
    \setcounter{enumi}{24}
    \item \ar{كعب يساوي أموالاً وجذوراً وأعداداً} --- Cube equals squares, roots, and number:
    $$x^3 = b x^2 + c x + d$$
\end{enumerate}
\end{propositionbox}

\clearpage
%% ========================================================================
\chapter{On the Geometric Solution of Cubic Equations}
\label{ch:geometric}
%% ========================================================================

%% ========================================================================
\section{The Fundamental Lemma}
\label{sec:lemma}
%% ========================================================================

Before presenting his solutions, Khayyam proves a fundamental lemma about solid figures (parallelopipeds):

\begin{propositionbox}[Lemma on Equal Solids \ar{(مقدمة في المجسمات المتساوية)}]

\bilingualline{%
إذا كان لنا مكعب مستطيل السطوح، قاعدته مربعة، وارتفاعه خط معلوم، وأردنا أن نبني على قاعدة مربعة أخرى مكعباً مستطيلاً مساوياً له، فنحن نجد خطاً يكون النسبة بين ضلع القاعدة الأولى وضلع القاعدة الثانية كالنسبة بين ذلك الخط والارتفاع.%
}{%
If we have a rectangular parallelopiped whose base is a square and whose height is a known line, and we wish to construct on another square base a rectangular parallelopiped equal to it, then we find a line such that the ratio between the side of the first base and the side of the second base is as the ratio between that line and the height.%
}

\end{propositionbox}

%% ========================================================================
\section{Solution of ``Cube and Roots Equal a Number'' ($x^3 + cx = d$)}
\label{sec:classic-case}
%% ========================================================================

This is the most celebrated case in Khayyam's treatise --- the equation $x^3 + c x = d$ (\ar{كعب وجذور يساويان عدداً}). Khayyam solves it by the intersection of a \textbf{parabola} and a \textbf{semicircle}.

\begin{propositionbox}[Case: $x^3 + c x = d$]

\bilingualline{%
فلنضرب مثالاً: كعب وجذور يساويان عدداً. فلنجعل العدد مكعباً مستطيل السطوح، قاعدته مربعة، أضلاعها مساوية لعدد الجذور، وارتفاعه خط معلوم.%
}{%
Let us give an example: cube and roots equal a number. Let us make the number a rectangular parallelopiped whose base is a square, whose sides are equal to the number of roots, and whose height is a known line.%
}

\vspace{0.3cm}

\bilingualline{%
ثم نأخذ قطعاً مكافئاً رأسه نقطة، ومحوره خط مستقيم، ومعلمه مساوٍ لضلع القاعدة. ونصف دائرة على ارتفاع المكعب، فنلقي من نقطة التقاطع القطع المكافئ والدائرة خطاً إلى المحور، فذلك الخط هو الجذر المطلوب.%
}{%
Then we take a parabola whose vertex is a point, whose axis is a straight line, and whose parameter is equal to the side of the base; and we describe a semicircle on the height of the parallelopiped. From the point of intersection of the parabola and the circle we drop a line to the axis: that line is the required root.\translatornote{The ``parameter'' (\textit{muʿallim}) of the parabola is its latus rectum. Khayyam's construction is equivalent to the modern method: setting $y^2 = bx$ (parabola) and $x^2 = y(h-y)$ (circle), whose intersection gives the solution.}%
}

\end{propositionbox}

%% ========================================================================
\subsection{The Construction}
\label{subsec:construction}
%% ========================================================================

The construction proceeds as follows:

\begin{enumerate}[noitemsep]
    \item Let $b^2 = c$ (the coefficient of $x$), and let $h$ be such that $b^2 h = d$ (the constant term)
    \item Construct a \textbf{parabola} with vertex $B$, axis $BZ$, and parameter $b$, so that for any point on the parabola, $y^2 = b x$
    \item Erect perpendicular $h$ at $B$, and on $h$ as diameter describe a \textbf{semicircle}
    \item Let the semicircle cut the parabola at point $D$
    \item From $D$ drop $DE$ perpendicular to the axis, and $DZ$ perpendicular to $BZ$
    \item Then $y = BE$ is the required root
\end{enumerate}

\begin{figure}[h]
\centering
\begin{tikzpicture}[scale=1.2,>=Stealth]
    % Axes
    \draw[->] (-0.3,0) -- (5,0) node[right] {$x$};
    \draw[->] (0,-0.3) -- (0,4.5) node[above] {$y$};
    
    % Parabola y^2 = 2x (parameter b=2)
    \draw[thick,domain=0:4,samples=100,variable=\x,blue] 
        plot ({\x},{sqrt(2*\x)}) node[right] {\small parabola};
    \draw[thick,domain=0:4,samples=100,variable=\x,blue] 
        plot ({\x},{-sqrt(2*\x)});
    
    % Semicircle on height h=3 at B
    \draw[thick,red] (0,0) arc (-90:90:1.5) node[above right] {\small semicircle};
    
    % Intersection point D (approximate)
    \coordinate (D) at (1.35, 1.65);
    \fill[black] (D) circle (1.5pt) node[above right] {$D$};
    
    % Drop perpendiculars
    \draw[dashed] (D) -- (1.35,0) node[below] {$Z$};
    \draw[dashed] (D) -- (0,1.65) node[left] {$E$};
    
    % Labels
    \node[below left] at (0,0) {$B$};
    \draw[<->] (5.3,0) -- (5.3,1.65) node[midway,right] {$y$};
    \draw[<->] (0,-0.6) -- (1.35,-0.6) node[midway,below] {$x$};
    
    % Parameter
    \node at (3.5,3.5) {$y^2 = bx$};
    \node at (2.5,0.8) {$x^2 = y(h-y)$};
\end{tikzpicture}
\caption{Khayyam's construction for $x^3 + cx = d$ by intersection of parabola and semicircle}
\label{fig:khayyam1}
\end{figure}

%% ========================================================================
\subsection{The Proof}
\label{subsec:proof}
%% ========================================================================

By the property of the parabola:
\[y^2 = bx \quad \Longrightarrow \quad \frac{x}{y} = \frac{y}{b}\]

By the property of the circle (with diameter $h$):
\[x^2 = y(h - y)\]

From these:
\[\frac{y}{b} = \frac{x}{y} = \frac{h - y}{x}\]

Therefore $b : y = y : x = x : (h - y)$ are in continued proportion.

From $b^2 : y^2 = b : (h - y)$ we get:
\[b^2(h - y) = y^3\]
\[b^2 h - b^2 y = y^3\]
\[y^3 + b^2 y = b^2 h\]

Substituting $c = b^2$ and $d = b^2 h$:
\[\boxed{y^3 + cy = d}\]

\clearpage
%% ========================================================================
\chapter{The Remaining Irreducible Cases}
\label{ch:remaining}
%% ========================================================================

%% ========================================================================
\section{Summary of Conic Section Pairs for Each Case}
\label{sec:summary}
%% ========================================================================

The following table summarizes the pairs of conic sections used for each of the 14 irreducible cases, organized according to Khayyam's five geometric families:

\begin{center}
\renewcommand{\arraystretch}{1.4}
\begin{longtable}{p{5.5cm}p{3.5cm}p{4cm}}
\toprule
\textbf{Equation} & \textbf{Conic Sections} & \textbf{Khayyam's Arabic Term} \\
\midrule
\endfirsthead
\toprule
\textbf{Equation} & \textbf{Conic Sections} & \textbf{Khayyam's Arabic Term} \\
\midrule
\endhead
\bottomrule
\endfoot
$x^3 + b^2 x = b^2 h$ & Circle + Parabola & \ar{دائرة ومكافئ} \\
$x^3 + b^2 h = b^2 x$ & Parabola + Hyperbola & \ar{مكافئ وزائد} \\
$x^3 = b^2 x + b^2 h$ & Parabola + Hyperbola & \ar{مكافئ وزائد} \\
$x^3 + ax^2 = c^3$ & Circle + Hyperbola & \ar{دائرة وزائد} \\
$x^3 + c^3 = ax^2$ & Circle + Hyperbola & \ar{دائرة وزائد} \\
$x^3 = ax^2 + c^3$ & Circle + Hyperbola & \ar{دائرة وزائد} \\
$x^3 + cx = ax^2 + d$ & Circle + Hyperbola & \ar{دائرة وزائد} \\
$x^3 + ax^2 = cx + d$ & Hyperbola + Hyperbola & \ar{زائدان} \\
$x^3 + d = ax^2 + cx$ & Circle + Hyperbola & \ar{دائرة وزائد} \\
$x^3 = ax^2 + cx + d$ & Hyperbola + Hyperbola & \ar{زائدان} \\
\bottomrule
\end{longtable}
\end{center}

\textit{Note:} \ar{مكافئ} = parabola (lit.\ ``equal'' or ``similar'', referring to the equality of abscissa and ordinate in the defining property); \ar{زائد} = hyperbola (lit.\ ``excess'', referring to the difference of squares in its defining property).

%% ========================================================================
\section{On the Existence and Multiplicity of Solutions}
\label{sec:existence}
%% ========================================================================

Khayyam carefully discusses conditions under which solutions exist. He observes that some cubic equations may have:
\begin{itemize}[noitemsep]
    \item No (positive real) solution
    \item Exactly one positive solution
    \item Two positive solutions
\end{itemize}

\begin{propositionbox}[On the Number of Solutions \ar{(في عدد الحلول)}]

\bilingualline{%
واعلم أن بعض هذه الأصناف قد يكون له حل واحد، وبعضها قد يكون له حلان، وبعضها لا يمكن أن يوجد إلا بشرط.%
}{%
Know that some of these species may have one solution, some may have two solutions, and some cannot be found except under a condition.%
}

\vspace{0.3cm}

\bilingualline{%
فإذا كان القطعان لا يلتقيان فلا حل للمسألة، وإذا كانا يلتقيان في نقطة واحدة فحل واحد، وإذا كانا يلتقيان في نقطتين فحلان.%
}{%
If the two conics do not intersect, there is no solution to the problem; if they intersect at one point, there is one solution; and if they intersect at two points, there are two solutions.\translatornote{Khayyam here anticipates the modern understanding that the number of real roots of a cubic corresponds to the number of intersection points of the two conics. He does not, however, have the concept of three real roots, as he restricts himself to positive solutions.}%
}

\end{propositionbox}

%% ========================================================================
\section{Case: Cube Equals Squares, Roots, and Number}
\label{sec:case25}
%% ========================================================================

The final case ($x^3 = b x^2 + c x + d$) is the most complex, requiring the intersection of two different hyperbolas:

\begin{propositionbox}[The Twenty-Fifth Species \ar{(الصنف الخامس والعشرون)}]

\bilingualline{%
وأما الصنف الخامس والعشرون: كعب يساوي أموالاً وجذوراً وأعداداً، فإنه يحتاج إلى قطعين زائدين مختلفين.%
}{%
As for the twenty-fifth species: cube equals squares, roots, and numbers, it requires two different hyperbolas.%
}

\vspace{0.3cm}

\bilingualline{%
فنصنع قطعاً زائداً له أحد محوريه على امتداد الخط المعلوم، ونصنع قطعاً زائداً آخر له أحد محوريه على امتداد الخط الآخر، ونلقي من نقطة التقاطعما خطاً إلى الموضع المعلوم، فهو الجذر المطلوب.%
}{%
We make a hyperbola one of whose axes is on the extension of the known line, and we make another hyperbola one of whose axes is on the extension of the other line; and from their point of intersection we drop a line to the known position, and it is the required root.%
}

\end{propositionbox}

\clearpage
%% ========================================================================
\chapter{Sample Problems and Applications}
\label{ch:problems}
%% ========================================================================

%% ========================================================================
\section{The Division of Ten into Two Parts}
\label{sec:division-ten}
%% ========================================================================

Khayyam illustrates his method with a problem leading to a cubic equation:

\begin{propositionbox}[Problem \ar{(مسألة)}]

\bilingualline{%
مسألة: قال قائل: قسّم عشرة إلى جزأين، بحيث يكون مجموع مربعي الجزأين مع خارج قسمة الكبير على الصغير اثنين وسبعين.%
}{%
A man said: Divide ten into two parts such that the sum of the squares of the two parts together with the quotient of the division of the greater by the smaller be seventy-two.%
}

\end{propositionbox}

\textit{Solution.} Let the greater part be $x$. Then the smaller is $10 - x$. The condition gives:
\[x^2 + (10 - x)^2 + \frac{x}{10 - x} = 72\]

Simplifying:
\[2x^2 - 20x + 100 + \frac{x}{10 - x} = 72\]

This leads to:
\[x^3 + 200x = 20x^2 + 2000\]

which is of the form $x^3 + cx = b x^2 + d$, solved by the intersection of a circle and a hyperbola.

%% ========================================================================
\section{On Doubling the Cube}
\label{sec:delian}
%% ========================================================================

Khayyam shows that the ancient Delian problem (doubling the cube) is equivalent to solving $x^3 = 2a^3$, which falls under his classification as ``cube equals number'':

\bilingualline{%
ومن ذلك مضاعفة المكعب، وهي مسألة مشهورة بين القدماء. فإنها تؤول إلى أن كعباً يساوي عدداً، والعدد هو مضاعفة المكعب المعطى.%
}{%
And among these is the duplication of the cube, which is a well-known problem among the ancients. It reduces to a cube equaling a number, where the number is a multiple of the given cube.\translatornote{The problem of doubling the cube (the ``Delian problem'') was one of the three famous classical Greek problems. Khayyam shows it reduces to $x^3 = 2a^3$, which is his Type 9 (cube equals number).}%
}

\clearpage
%% ========================================================================
\chapter{Concluding Remarks}
\label{ch:conclusion}
%% ========================================================================

Khayyam concludes his treatise with reflections on the nature of algebra and its relationship to geometry:

\begin{propositionbox}[On the Relation of Algebra to Geometry \ar{(في علاقة الجبر بالهندسة)}]

\bilingualline{%
وقد أتممنا ما أردنا من جمع أصناف المسائل الجبرية، ووضع البراهين عليها.%
}{%
We have completed what we intended of gathering the species of algebraic problems and setting down proofs for them.%
}

\vspace{0.3cm}

\bilingualline{%
ونحن نعلم أن بعض الناس يعتقدون أن الجبر بريء من البرهان الهندسي، ولكن هذا باطل.%
}{%
We know that some people believe that algebra is free from geometrical proof, but this is false.%
}

\vspace{0.3cm}

\bilingualline{%
فإن الأصناف التي ذكرناها لا يمكن أن تُبرهن إلا بالقطوع المخروطية أو بغيرها من الأدلة الهندسية.%
}{%
For the species we have mentioned cannot be proved except by conic sections or by other geometrical arguments.%
}

\vspace{0.3cm}

\bilingualline{%
والجبر والمقابلة إنما هما وسيلة لاستخراج المجهول، والبرهان على صحة ذلك إنما هو بالهندسة.%
}{%
Algebra and Almucabala are but a means of extracting the unknown, and the proof of its correctness is by geometry.\translatornote{This passage is one of the most celebrated in Khayyam's treatise. It reveals his deep understanding that algebraic manipulation alone is insufficient without geometric proof --- a view that would not be fully superseded until the development of symbolic algebra in the 16th--17th centuries.}%
}

\end{propositionbox}

\vspace{1cm}

\begin{center}
\bilingualline{%
والله ولي التوفيق.%
}{%
And God is the source of all success.%
}
\end{center}

\vspace{1cm}

\begin{center}
{\color{rulecolor}\rule{0.6\textwidth}{0.8pt}}

\vspace{0.5cm}

{\small\textit{End of the Treatise}}
\end{center}

\clearpage
%% ========================================================================
\chapter*{Editor's and Translator's Notes}
\addcontentsline{toc}{chapter}{Editor's and Translator's Notes}
%% ========================================================================

This bilingual edition is based primarily on the 13th-century Lahore manuscript (Smith Oriental MS 45, Columbia University Rare Book and Manuscript Library), with reference to the following sources:

\begin{itemize}[noitemsep]
    \item Woepcke, F. \textit{L'Alg\`{e}bre d'Omar Alkhayy\^{a}m\^{i}}. Paris: B. Duprat, 1851.
    \item Kasir, D. S. \textit{The Algebra of Omar Khayyam}. New York: Teachers College, Columbia University, 1931.
    \item Rashed, R. and Vahabzadeh, B. \textit{Omar Khayyam, the Mathematician}. New York: Bibliotheca Persica Press, 2000.
\end{itemize}

\section*{Translation Methodology}

The English translation follows the style of D. S. Kasir's 1931 edition, which remains the standard English version. Key translation choices:

\begin{itemize}
    \item \textbf{\ar{كعب} (kaʿb):} Rendered as ``cube'', following the literal meaning of the Arabic (a solid with six square faces, like a gaming die). Khayyam uses this for $x^3$.
    
    \item \textbf{\ar{مال} (māl):} Rendered as ``square''. The literal meaning is ``wealth'', ``property'', or ``possession''. In Arabic algebra, the square of an unknown was metaphorically termed ``possession'' because it represents the amount owned.
    
    \item \textbf{\ar{جذر} (jadhr):} Rendered as ``root''. The word literally means ``root'' (of a plant) or ``base'' or ``origin''. In Arabic algebra, $x$ is the root from which the square and cube grow.
    
    \item \textbf{\ar{زيادة} (ziyāda):} Rendered as ``augmentation'' or ``excess''. In Khayyam's geometric terminology, the hyperbola (\ar{زائد}) is named for the ``excess'' property in its defining proportion.
    
    \item \textbf{\ar{ناقص} (nāqiṣ):} Rendered as ``diminution'' or ``defect''. In quadratic equations, this term denotes the subtracted quantity.
\end{itemize}

\section*{On Khayyam's Geometric Algebra}

Khayyam's treatise represents the culmination of the Greek-Arabic tradition of \textit{geometric algebra}. Unlike modern symbolic algebra, Khayyam's methods are fundamentally geometric: every algebraic quantity is represented as a geometric magnitude (line, surface, or solid), and every operation is justified by a geometric proof.

The key insight of the treatise is that cubic equations of the form $x^3 + ax = b$ (and the other 13 irreducible types) \textit{cannot} be solved by ruler and compass alone. Khayyam proves that the intersection of conic sections is necessary and sufficient for their solution --- a result that would not be fully appreciated in Europe until the 17th century.

\vspace{1cm}

\begin{center}
{\color{rulecolor}\rule{0.4\textwidth}{0.4pt}}

\vspace{0.5cm}

{\footnotesize\textit{Bilingual edition typeset in \LaTeX\ using XeLaTeX, polyglossia, and Noto Naskh Arabic}}

{\footnotesize\textit{Mathematical typesetting via amsmath and TikZ}}

{\footnotesize\textit{Arabic text normalized for readability while preserving technical vocabulary}}
\end{center}

\end{document}
%% ========================================================================

